\documentclass{article}

\usepackage{amssymb}
\usepackage{amsthm}
\usepackage{amsmath}
\usepackage{geometry}
\usepackage{mathrsfs}
\usepackage{stmaryrd}
\usepackage{tikz-cd}

\usepackage[linktoc=all, hidelinks]{hyperref}

\newcommand{\CC}{\mathbb{C}}
\newcommand{\FF}{\mathbb{F}}
\newcommand{\NN}{\mathbb{N}}
\newcommand{\QQ}{\mathbb{Q}}
\newcommand{\RR}{\mathbb{R}}
\newcommand{\ZZ}{\mathbb{Z}}
\newcommand{\id}{\mathrm{id}}
\newcommand{\op}[1]{#1^{\text{op}}}
\newcommand{\powerset}[1]{\mathcal{P}(#1)}
\newcommand{\tensor}{\otimes}
\newcommand{\rad}[1]{\sqrt{#1}}

\DeclareMathOperator{\Ann}{Ann}
\DeclareMathOperator{\Spec}{Spec}
\DeclareMathOperator{\Max}{Max}

% ideal shortcuts
\newcommand{\ideala}{\mathfrak{a}}
\newcommand{\idealb}{\mathfrak{b}}
\newcommand{\idealc}{\mathfrak{c}}
\newcommand{\idealm}{\mathfrak{m}}
\newcommand{\idealp}{\mathfrak{p}}
\newcommand{\idealq}{\mathfrak{q}}
\newcommand{\nilrad}{\mathfrak{R}}
\newcommand{\jacobrad}{\mathfrak{R}}

\DeclareMathOperator{\res}{res}

\newcommand{\Set}{\mathbf{Set}}
\newcommand{\Ab}{\mathbf{Ab}}
\newcommand{\Grp}{\mathbf{Grp}}
\newcommand{\Mod}{\mathbf{Mod}}

\renewcommand\labelenumii{(\roman{enumii})}

\theoremstyle{definition}
\newtheorem{theorem}{Theorem}[section]
\newtheorem{proposition}[theorem]{Proposition}
\newtheorem{definition}[theorem]{Definition}
\newtheorem{lemma}[theorem]{Lemma}
\newtheorem{corollary}[theorem]{Corollary}
\newtheorem{example}{Example}[theorem]

\newtheorem{exercise}{Exercise}[section]


\title{Atiyah Macdonald exercises}
\author{bpaul}
\date{\today}

\begin{document}
\maketitle

\section{Rings and Ideals}

\setcounter{exercise}{11}
\renewcommand{\theexercise}{\thesection.\arabic{exercise}}

% 1.12
\begin{exercise}
    \begin{enumerate}
        \item \(\ideala \subseteq (\ideala : \idealb)\)
        \item \((\ideala : \idealb) \idealb \subseteq \ideala\)
        \item \(((\ideala : \idealb) : \idealc) = (\ideala : \idealb \idealc) = ((\ideala : \idealc) : \idealb)\)
        \item \((\bigcap_i \ideala_i : \idealb) = \bigcap_i (\ideala_i : \idealb)\)
        \item \((\ideala : \sum_i \idealb_i) = \bigcap_i (\ideala : \idealb_i)\).
    \end{enumerate}
\end{exercise}

\noindent\textbf{Solution:}
\begin{enumerate}
    \item
        For \(a \in \ideala\), \(a \idealb \subseteq \ideala\) since \(\ideala\) is an ideal.
    \item
        Since \(xb \in (\ideala : \idealb) \idealb\) iff \(x \idealb \subseteq \ideala\) and \(b \in \idealb\), \(xb \in \ideala\).
    \item
        Since \(\idealb \idealc = \idealc \idealb\) we only need to show the first equality.
        Observe that \(x \in ((\ideala : \idealb) : \idealc)\) iff \(x \idealc \subseteq (\ideala : \idealb)\), that is for all \(xc \in x \idealc\) with \(c \in \idealc\), \(xc \idealb \subseteq \ideala\).
        This is true iff \(x \idealc \idealb \subseteq \ideala\).
    \item
        \(x \idealb \subseteq \bigcap_i \ideala_i\) iff \(x \idealb \subseteq \ideala_i\) for all \(i\).
    \item
        Note \(x \sum_i \idealb_i = \sum_i x \idealb_i\).
        Then \(\sum_i x \idealb_i \subseteq \ideala\) iff \(x \idealb_i \subseteq \ideala\) for all \(i\).
\end{enumerate}

% 1.13
\begin{exercise}
    \begin{enumerate}
        \item \label{ex-1.13.1} \(\rad{\ideala} \supseteq \ideala\)
        \item \(\rad{\rad{\ideala}} = \rad{\ideala}\)
        \item \(\rad{\mathfrak{ab}} = \rad{\ideala \cap \idealb} = \rad{\ideala} \cap \rad{\idealb}\)
        \item \(\rad{\ideala} = (1) \iff \ideala = (1)\)
        \item \(\rad{\ideala + \idealb} = \rad{\rad{\ideala} + \rad{\idealb}}\)
        \item if \(\idealp\) is prime, \(\rad{\idealp^n} = \idealp\) for all \(n > 0\).
    \end{enumerate}
\end{exercise}

\noindent\textbf{Solution:}
\begin{enumerate}
    \item
        If \(x \in \ideala\), then \(x^1 \in \ideala\) so \(x \in \rad{\ideala}\).
    \item
        \(x \in \rad{\rad{\ideala}}\) iff \((x^n)^m \in \ideala\) iff \(x^{nm} \in \ideala\) iff \(x \in \rad{\ideala}\).
    \item
        If \(x^n \in \ideala \idealb\) then \(x^n \in \ideala \cap \idealb\).
        If \(x^n \in \ideala \cap \idealb\) then \(x^n \in \ideala\) and \(x^n \in \idealb\).
        If \(x^n \in \ideala\) and \(x^m \in \idealb\) then \(x^{mn} \in \ideala \idealb\).
    \item
        If \(1^n \in \ideala\) then \(1 \in \ideala\).
        The converse comes from \ref{ex-1.13.1}.
    \item
        By 1., \(\rad{\ideala + \idealb} \subseteq \rad{\rad{\ideala} + \rad{\idealb}}\).
        If \(x^n = y + z\) where \(y^k \in \ideala\) and \(z^m \in \idealb\) then \(x^{n (k + m)} = (y + z)^{k + m}\) has every term in its binomial expansion with either its \(y\) term at least \(k\) or its \(z\) term at least \(m\), and so \(x^{n (k+m)} \in \ideala + \idealb\).
    \item
        If \(x^m \in \idealp^n\) then \(x^m \in \idealp\) so \(x \in \idealp\).
        Conversely if \(x \in \idealp\) then \(x^n \in \idealp^n\) so \(x \in \rad{\idealp^n}\).
\end{enumerate}

\setcounter{exercise}{17}

% 1.18
\begin{exercise}
    If \(\ideala_1, \ideala_2\) are ideals of \(A\) and if \(\idealb_1, \idealb_2\) are ideals of \(B\), then
    \[
        \begin{array}{c c}
            (\ideala_1 + \ideala_2)^e = \ideala_1^e + \ideala_2^e, &
            (\idealb_1 + \idealb_2)^c \supseteq \idealb_1^c + \idealb_2^c, \\
            (\ideala_1 \cap \ideala_2)^e \subseteq \ideala_1^e \cap \ideala_2^e, &
            (\idealb_1 \cap \idealb_2)^c = \idealb_1^c \cap \idealb_2^c, \\
            (\ideala_1 \ideala_2)^e = \ideala_1^e \ideala_2^e, &
            (\idealb_1 \idealb_2)^c \supseteq \idealb_1^c \idealb_2^c, \\
            (\ideala_1 : \ideala_2)^e \subseteq (\ideala_1^e : \ideala_2^e), &
            (\idealb_1 : \idealb_2)^c \subseteq (\idealb_1^c : \idealb_2^c), \\
            \rad(\ideala)^e \subseteq \rad(\ideala^e) &
            \rad(\idealb)^c = \rad(\idealb^c).
        \end{array}
    \]
\end{exercise}

\noindent\textbf{Solution:}
I REALLY DON'T WANT TO!!!
remind me to do this later please

\setcounter{exercise}{0}
\renewcommand{\theexercise}{(\arabic{exercise})}

\begin{exercise} \label{ex1-1}
    Let \(x\) be a nilpotent element of a ring \(A\).
    Show that \(1 + x\) is a unit in \(A\).
    Deduce that the sum of a nilpotent element and a unit is a unit.
\end{exercise}

\noindent\textbf{Solution:}
Suppose \(x^n = 0\).
Notice that
\[
    (1 + x)(\prod_{i=0}^n (-1)^i x^i) = 1 + x - x - x^2 + x^2 + x^3 \cdots \pm x^n = 1 \pm x^n = 1.
\]
If \(u\) is a unit then \(u^{-1} x\) is nilpotent, so \(1 + u^{-1} x\) is a unit, hence \(u + x\) is.

\begin{exercise}
    Let \(A\) be a ring and let \(A[x]\) be the ring of polynomials in an indeterminate \(x\), with coefficients in \(A\).
    Let \(f = a_0 + a_1 x + \cdots + a_n x^n\). Prove that
    \begin{enumerate}
        \item \label{ex1-2-1}
            \(f\) is a unit \(\iff\) \(a_0\) is a unit in \(A\) and \(a_1, \cdots, a_n\) are nilpotent.
        \item
            \(f\) is nilpotent \(\iff\) \(a_0, a_1, \cdots, a_n\) are nilpotent.
        \item
            \(f\) is a zero-divisor \(\iff\) there exists \(a \neq 0\) in \(A\) such that \(a f = 0\).
            [Choose a polynomial \(g = b_0 + b_1 x + \cdots + b_m x^m\) of least degree \(m\) such that \(f g = 0\).
            Then \(a_n b_m = 0\), hence \(a_n g = 0\) because \(a_n g\) annihilates \(f\) and has degree \(< m\).
            Now show by induction that \(a_{n-r} g = 0 \, (0 \leq r \leq n)\).]
        \item
            \(f\) is said to be primitive if \((a_0, a_1, \cdots, a_n) = (1)\).
            Prove that if \(f, g \in A[x]\), then \(fg\) is primitive \(\iff\) f and g are primitive.
    \end{enumerate}
\end{exercise}

\noindent\textbf{Solution:}
\begin{enumerate}
    \item
        If \(a_0\) is a unit and \(a_i\) are nilpotent for \(i \geq 1\), then \(a_i x^i\) are nilpotent and their sum is nilpotent, hence by \ref{ex1-1} \(f\) is a unit.

        Conversely suppose \(f\) is a unit with inverse \(g\) (coefficients indexed with \(b_i\)).
        Then \(a_0 b_0 = 1\).
        Next let \(\idealp\) be a prime ideal.
        The map \(\varphi : A[x] \to (A/\idealp)[x]\) has \(\varphi(f)\varphi(g) = 1\).
        Hence the degree of \(\varphi(f)\) and \(\varphi(g)\) are both \(0\) since \((A/\idealp)[x]\) is an integral domain.
        Thus \(\phi(a_i) = 0\) for all \(i \geq 1\), that is \(a_i \in \idealp\).
        Since this is for any prime ideal, the \(a_i\) are in all prime ideals, and hence are nilpotent.
    \item
        Similarly to above, if the \(a_i\) are nilpotent, their product with indeterminates are, and their sum is.

        Let \(f\) be nilpotent with \(f^n = 0\).
        Then \(a_0^n = 0\) so \(a_0\) is nilpotent.
        By exercise \ref{ex1-1}, \(f + 1\) is a unit, and by part \ref{ex1-2-1} the \(a_i\) are nilpotent for \(i \geq 1\).
    \item
        Let \(g \neq 0\) annihilate \(f\) with minimum degree \(m\).
        Then \(b_m a_n = 0\).
        Hence \(a_n g\) has degree less than m, and moreover \(a_n g\) annihilates \(f\).
        Hence \(a_n g = 0\).
        Now suppose for induction that \(a_{n-i} g = 0\) holds for all \(i \leq r\).
        Then \(a_{n-r-1} b_i = - \sum_{j > i} a_{n-j} b_j = 0\).
        Hence for all \(i\), \(a_i b_n = 0\).
        Thus \(b_n f = 0\).
    \item
        Suppose \(fg\) is primitive, and call its coefficients \(c_i\).
        It is sufficient to show \(f\) is primitive.
        Notice that \((a + b) \subseteq (a) + (b)\).
        Hence \((c_i) \subseteq \sum_{j \leq i} (a_j b_{i-j}) \subseteq \sum_{j \leq i} (a_j)\) for all \(i\).
        Thus \((1) = \sum_i (c_i) \subseteq \sum_i (a_j)\), and so \(f\) is primitive.

        Conversely suppose \(f\) and \(g\) are primitive.
        We will show that \(\idealc = (c_0, c_1, \cdots, c_k)\) is not contained in any prime ideal of \(A\).
        Suppose it were contained \(\idealp\).
        Then under the map \(\varphi : A[x] \to (A/\idealp)[x]\), \(\varphi(c_i) = 0\) and so \(\varphi(fg) = 0\).
        Since \((A/\idealp)[x]\) is a domain, either \(\varphi(f) = 0\) or \(\varphi(g) = 0\).
        However, since \((a_0, a_1, \cdots, a_n) \not\supseteq \idealp\) (since \(\idealp\) is proper), some \(a_i\) is not sent to zero.
        Similarly, some \(b_i\) is not sent to zero.
        This is a contradiction, and so \(\idealc\) is not contained in any prime ideal.
        Since every proper ideal is contained in a maximal (and hence prime) ideal, \(\idealc\) is not proper.
\end{enumerate}

\section{Modules}
\renewcommand{\theexercise}{\thesection.\arabic{exercise}}

\end{document}
